\begin{proof}[Proof of \cref{obj:lem:empirical-threshold-process-bounded}]
Fix \(n,d,\epsilon\) and the sample. For each cell \(j\), \cref{obj:def:dual-process,obj:def:jf-estimator} give, for \(\lambda\in[0,1]\),
\[
f_\lambda(\widehat q_j)
=
g_0(\widehat q_j)
+
\max\!\left\{0,\,
g_1(\widehat q_j)-g_0(\widehat q_j)
-\lambda s(\widehat q_j)\right\}.
\]
The empirical vector \(\widehat q_j\) is fixed as \(\lambda\) varies. Thus this expression is continuous in \(\lambda\), and so is its finite sum \(\widehat F(\lambda)\). When \(n=0\), the empirical vectors are zero under the convention \(x/0=0\); when \(d=0\), the sum is empty and \(\widehat F\) is identically zero. The same continuity argument covers both cases. % lean: CausalSmith.Stat.DiscreteBudgetvalueCurve.thresholdFunReal_continuous_time

Since \([0,1]\) is compact, the continuous real-valued function \(\widehat F\) is bounded there. Hence some \(B\in\mathbb R\) satisfies
\[
|\widehat F(\lambda)|\leq B
\qquad\text{for every }\lambda\in[0,1].
\]
% lean: CausalSmith.Stat.DiscreteBudgetvalueCurve.empiricalThresholdProcess_bounded
\end{proof}
